\chapter{Where to Begin}
\label{ch:orientation}

\chapterguide%
  {No previous machine-learning knowledge is required.}
  {identify observations, features, labels, models, metrics, learning settings, and the project workflow.}

Machine learning builds systems that learn useful patterns from examples rather than requiring a rule for every situation.

\begin{intuition}
Learn house prices from examples of sold homes rather than writing a rule for every possible house.
\end{intuition}

\section{The five parts of a machine-learning problem}

\begin{mltable}
\begin{tabularx}{\linewidth}{@{}>{\bfseries\leavevmode\color{MLNavy}}p{0.20\linewidth}X@{}}
\toprule
\rowcolor{MLPaleBlue!92}
Part & Meaning \\
\midrule
Observation & One example or row, such as one house, patient, customer, or email. \\
Feature & An input used to make a prediction, such as floor area, age, or word count. \\
Target & The answer the model should learn to predict, such as price or spam/not spam. \\
Model & A mathematical rule that turns features into a prediction. \\
Metric & A number used to judge whether the model is useful on unseen data. \\
\bottomrule
\end{tabularx}
\end{mltable}

\begin{mltable}
\begin{tabularx}{\linewidth}{@{}>{\bfseries\leavevmode\color{MLNavy}}p{0.14\linewidth}X X X >{\bfseries}X@{}}
\toprule
\rowcolor{MLPaleBlue!92}
House & Area & Bedrooms & Distance to center & Sale price \\
\midrule
A & 90 m$^2$ & 2 & 8 km & \textdollar 220,000 \\
B & 140 m$^2$ & 3 & 4 km & \textdollar 360,000 \\
C & 110 m$^2$ & 2 & 12 km & \textdollar 245,000 \\
\bottomrule
\end{tabularx}
\end{mltable}

Each house is an \term{observation}; the middle columns are \term{features} known before sale, and price is the \term{target}. Training learns from completed rows; prediction estimates a new row's missing price.

\begin{workedexample}
Predict sale price from floor area, bedrooms, and distance to the center; evaluate with mean absolute error.
\end{workedexample}

\section{Parameters, hyperparameters, loss, and metrics}

\begin{mltable}
\begin{tabularx}{\linewidth}{@{}>{\bfseries\leavevmode\color{MLNavy}}p{0.2\linewidth}X X@{}}
\toprule
\rowcolor{MLPaleBlue!92}
Term & Who chooses it? & Example \cr
\midrule
Parameter & Learned from training data & A regression coefficient or a tree split \cr
Hyperparameter & Chosen before or during model selection & Number of neighbors, tree depth, regularization strength \cr
Loss & Used by the training process to measure wrongness & Squared error, cross-entropy \cr
Metric & Used to report practical performance & Accuracy, recall, MAE, RMSE \cr
\bottomrule
\end{tabularx}
\end{mltable}

Training loss guides fitting; the reporting metric measures practical success. They may differ.

\section{The complete workflow}

\begin{mlfigure}
\begin{tikzpicture}[
  node distance=0.45cm and 0.55cm,
  box/.style={draw=MLBlue,fill=MLPaleBlue,rounded corners=2mm,text width=3.35cm,minimum height=0.9cm,inner xsep=3pt,align=center,font=\scriptsize\bfseries\color{MLNavy}},
  arr/.style={-{Stealth[length=2.2mm]},thick,draw=MLTeal}
]
\node[box] (a) {1. Define the problem};
\node[box,right=of a] (b) {2. Understand the data};
\node[box,right=of b] (c) {3. Split correctly};
\node[box,right=of c] (d) {4. Prepare the features};
\node[box,below=0.65cm of d] (e) {5. Build a baseline};
\node[box,left=of e] (f) {6. Train and tune models};
\node[box,left=of f] (g) {7. Evaluate fairly};
\node[box,left=of g] (h) {8. Explain and monitor};
\draw[arr] (a) -- (b); \draw[arr] (b) -- (c); \draw[arr] (c) -- (d);
\draw[arr] (d) -- (e); \draw[arr] (e) -- (f); \draw[arr] (f) -- (g); \draw[arr] (g) -- (h);
\end{tikzpicture}
\end{mlfigure}

\begin{mlsteps}
\setlength{\itemsep}{0.12em}
  \item \term{Define the decision.} State what must be predicted, who will use the result, and what a wrong prediction costs.
  \item \term{Validate the data.} Check rows, keys, units, missingness, timing, and bias. Apply only fixed, label-blind corrections before splitting.
  \item \term{Split the data.} Create training, validation, and test data in a way that matches future use.
  \item \term{Prepare the features.} Learn imputation, category encoding, scaling, and other data-dependent transformations from training data inside a pipeline.
  \item \term{Build a baseline.} Start with a simple rule, such as predicting the average or most common class.
  \item \term{Train and tune.} Fit candidate models on training data and choose settings with validation data.
  \item \term{Evaluate once on the test set.} Use the untouched test data for the final estimate.
  \item \term{Explain, deploy, and monitor.} Document limitations and check whether the data or performance changes over time.
\end{mlsteps}

Fit preprocessing on training data only; never select models using the test set (\secref{sec:data-splits}, \secref{sec:data-leakage}).

\section{The main learning settings}

\begin{mltable}
\small
\renewcommand{\arraystretch}{1.12}
\begin{tabularx}{\linewidth}{@{}>{\bfseries\leavevmode\color{MLNavy}}p{0.29\linewidth}X@{}}
\toprule
\rowcolor{MLPaleBlue!92}
Learning setting & How it works \\
\midrule
Supervised learning & Learns from labeled examples to predict values or classes. \\
Unsupervised learning & Finds groups, patterns, or unusual cases without labels. \\
Semi-supervised learning & Combines a small labeled set with a larger unlabeled set. \\
Reinforcement learning & Learns actions through interaction and rewards over time. \\
\bottomrule
\end{tabularx}
\end{mltable}

Supervised and unsupervised learning form the core; limited-label methods, reinforcement learning, and neural networks follow.

\begin{intuition}
\term{Self-supervised learning} creates targets from the data, such as hidden words or image regions, to learn representations. \term{Active learning} chooses which examples to send for human labeling. The field map places both in context.
\end{intuition}

\begin{mlfigure}
\begin{tikzpicture}[
  node distance=0.42cm and 0.55cm,
  hd/.style={draw=MLNavy,fill=MLNavy,text=white,rounded corners=1.5mm,minimum width=3.35cm,minimum height=0.62cm,align=center,font=\small\bfseries},
  bx/.style={draw=MLBlue,fill=MLPaleBlue,rounded corners=1.5mm,minimum width=3.35cm,minimum height=1.5cm,align=center,font=\scriptsize\color{MLNavy}},
  arr/.style={-{Stealth[length=2mm]},thick,draw=MLTeal}
]
\node[hd] (h1) {Supervised};
\node[hd,right=of h1] (h2) {Semi-supervised};
\node[hd,right=of h2] (h3) {Unsupervised};
\node[hd,right=of h3] (h4) {Reinforcement};
\node[bx,below=0.28cm of h1] (b1) {You are given the\\answer for every example.\\[0.25em]\textbf{``This email is spam.''}\\[0.25em]Goal: copy the mapping};
\node[bx,below=0.28cm of h2] (b2) {You are given answers\\for only some examples.\\[0.25em]\textbf{``A few emails are labeled.''}\\[0.25em]Goal: learn from both sets};
\node[bx,below=0.28cm of h3] (b3) {You are given no answers\\at all, only the examples.\\[0.25em]\textbf{``Here are 10\,000 emails.''}\\[0.25em]Goal: find the structure};
\node[bx,below=0.28cm of h4] (b4) {You are given a score\\after you act.\\[0.25em]\textbf{``That move lost you 3 points.''}\\[0.25em]Goal: act better next time};
\end{tikzpicture}
\end{mlfigure}

\section{Generalization: the real goal}

A model \term{generalizes} when it performs well on new examples from the intended setting. High training performance may reflect memorization instead.

\begin{intuition}
Memorizing practice answers does not prepare a student for new questions. Models must generalize too.
\end{intuition}

Three outcomes are common:

\begin{mltable}
\begin{tabularx}{\linewidth}{@{}>{\bfseries\leavevmode\color{MLNavy}}p{0.15\linewidth}>{\raggedright\arraybackslash}X@{}}
\toprule
\rowcolor{MLPaleBlue!92}
Outcome & Model behavior \\
\midrule
Underfitting & The model is too simple and performs poorly on both training and unseen data. \\
Good fit & The model captures useful structure and performs similarly on training and unseen data. \\
Overfitting & The model performs very well on training data but poorly on unseen data. \\
\bottomrule
\end{tabularx}
\end{mltable}

\section{What a wrong prediction costs}

State the cost of each error type before choosing a model.

\begin{mltable}
\begin{tabularx}{\linewidth}{@{}>{\bfseries\leavevmode\color{MLNavy}}p{0.25\linewidth}XX@{}}
\toprule
\rowcolor{MLPaleBlue!92}
\normalfont\bfseries\leavevmode\color{MLNavy}Problem & \normalfont\bfseries\leavevmode\color{MLNavy}Cost of a false alarm & \normalfont\bfseries\leavevmode\color{MLNavy}Cost of a miss \\
\midrule
Spam filter & A real email is hidden --- irritating and occasionally serious & A spam email arrives --- mildly annoying \\
Disease screening & An unnecessary follow-up test --- costly and stressful & An untreated illness --- potentially fatal \\
Fraud detection & A legitimate card is blocked --- an angry customer & Money is lost --- direct financial damage \\
Movie recommendation & A poor suggestion --- ignored & A good movie is not shown --- an invisible loss \\
\bottomrule
\end{tabularx}
\end{mltable}

\begin{intuition}
Error costs determine which metrics matter. \chapref{ch:evaluation} develops precision, recall, ranking, and thresholds for evaluating those trade-offs.
\end{intuition}

\begin{tryit}
Choose an automated decision. Identify its observation, features, target, model, metric, and the cost of each error type.
\end{tryit}

\section{What to learn first}

Follow the course map: tools and foundations first, each workflow before its models, then advanced topics and the capstone.

\begin{keyidea}
A useful model requires a well-defined problem and fair evaluation, not just an algorithm.
\end{keyidea}

\begin{quickcheck}
\begin{enumerate}
  \item What are the observation, features, target, model, and metric in a problem that predicts whether a loan payment will be late?
  \item Why is a learned model parameter different from a hyperparameter chosen during model selection?
  \item Why can excellent training performance still be evidence of a poor machine-learning system?
\end{enumerate}
\end{quickcheck}
